%Konfigurationsdatei für Angaben

% Titel, Autor
\setmytitle{Reduzierung redundanter Datenzugriffe in einem System zur Coupon-Validierung bei internationalen Reisen}
\setmyauthor{Friedrich Lidzba}
\setmyimmanr{s3005729}

% Studienrichtung
\setmyfos{Informationstechnik}
% Studiengang
\setmycos{3IT24-2}

% Datum der Themenübergabe
\setmytsdate{03.09.2026}
% Datum der Abgabe
\setmysdate{05.11.2026}
\setmyyear{\the\year} %für Links in der Fußnote

%Semester
\setmysemester{4}
%Gutachter (mit Titel!)
\setmyexpert{Christian Heidrich, M.SC.}

% Anschrift der Firma (Ort wird auch für Unterschriften verwendet)
\setmycompany{Deutsche Telekom AG}
\setmystreet{Riesaerstraße 5}
\setmyplz{01129}
\setmylocation{Dresden} % Stadt

% Ob Sperrvermerk, Abbildungsverzeichnis, Glossar, Abkürzungsverzeichnis verwendet werden(true heißt mit, false ohne)
\setbool{blockingNotice}{false} %Sperrvermerk
\setbool{lof}{true} %Abbildungsverzeichnis
\setbool{glossary}{true} %Glossar
\setbool{loacr}{true} %Abkürzungsverzeichnis

% Aussehen der Verlinkungen und festlegen des Titels
\hypersetup{
    pdftitle        = {\mytitle},
    pdfauthor       = {\myauthor},
    pdfsubject		= {Belegarbeit an der Dualen Hochschule Sachsen},
    pdfcenterwindow	= {true},
    colorlinks		= {false}, %Links farbig malen
    linkcolor		= {HKS44-100},
    citecolor		= {HKS57-100},
    filecolor		= {HKS07-100},
    urlcolor		= {HKS44-100},
}

%eigene Kurzbefehle -----------------------------------------
\newcommand{\az}[1]{\glqq{}#1\grqq{}} %Anführungszeichen

\newcommand{\dhsn}{\ac{dhsn}} %Schreibfaul, automatisch "DHSN" geschrieben